不是独立的或不是精确的卡方分布, 把它们当做如此来作图, 也是有用的. 所得到的图称为卡方图或伽马图, 因为卡方分布是更一般的伽马分布的特殊情况(参见文献[6]).

构造卡方图的步骤如下:

\begin{enumerate}[label=\arabic*.,leftmargin=2em,itemsep=2pt]
\item 把\hyperref[eq:4-32]{式(4-32)}中的平方距离从最小到最大排序为 $d_{(1)}^{2}\le d_{(2)}^{2}\le\cdots\le d_{(n)}^{2}$.
\item 画出点对 $\Bigl(q_{c,p}\Bigl(\dfrac{j-\frac{1}{2}}{n}\Bigr),d_{(j)}^{2}\Bigr)$ 的图形, 其中 $q_{c,p}\Bigl(\dfrac{j-\frac{1}{2}}{n}\Bigr)$ 是自由度为 $p$ 的卡方分布的 $100\Bigl(\dfrac{j-\frac{1}{2}}{n}\Bigr)$ 分位数.
\end{enumerate}

分位数是用比例来规定的, 而百分位数是用百分比来规定的.

分位数 $q_{c,p}\Bigl(\dfrac{j-\frac{1}{2}}{n}\Bigr)$ 与卡方分布的上百分位数有关. 具体的, $q_{c,p}\Bigl(\dfrac{j-\frac{1}{2}}{n}\Bigr)=\chi_{p}^{2}\Bigl(\dfrac{n-j+\frac{1}{2}}{n}\Bigr)$.

图形应当类似于一条通过原点且斜率为 1 的直线. 一个系统的弯曲图形表明缺乏正态性. 一个或两个远离直线的点表示大的距离, 或是离群的观测值, 这些都值得进一步注意.

\noindent\textbf{例 4.13}\quad \textbf{(构造卡方图)}

让我们构造在例 4.12 中所给的广义距离的卡方图. 对于 $p=2$ 和 $n=10$ 的有序距离和相应的卡方百分位数列表如下:

\begin{mvtable}
\centering
\mvfit{%
\begin{tabular}{cccccc}
\toprule
$j$ & $d_{(j)}^{2}$ & $q_{c,2}\Bigl(\dfrac{j-\frac{1}{2}}{10}\Bigr)$ & $j$ & $d_{(j)}^{2}$ & $q_{c,2}\Bigl(\dfrac{j-\frac{1}{2}}{10}\Bigr)$ \\
\midrule
1 & 0.30 & 0.10 & 6 & 1.64 & 1.60 \\
2 & 0.62 & 0.33 & 7 & 1.71 & 2.10 \\
3 & 1.16 & 0.58 & 8 & 1.79 & 2.77 \\
4 & 1.30 & 0.86 & 9 & 3.53 & 3.79 \\
5 & 1.61 & 1.20 & 10 & 4.38 & 5.99 \\
\bottomrule
\end{tabular}
}
\end{mvtable}

点对 $\Bigl(q_{c,2}\Bigl(\dfrac{j-\frac{1}{2}}{10}\Bigr),d_{(j)}^{2}\Bigr)$ 的图形示于\hyperref[fig:4.7]{图 4.7}. \hyperref[fig:4.7]{图 4.7} 中的点基本上在一条直线上, 但给定的样本容量小, 我们从该图形的特征很难来拒绝数据的二元正态性. 如果需要对数据作进一步分析, 把它们变换为更接近二元正态的观测值是合理的. 适当的变换方法在 4.8 节中讨论.

除了检验单变量图和散布图外, 我们应通过构造卡方图或 $d^{2}$ 图来检验多元正态性. \hyperref[fig:4.8]{图 4.8} 给出了对计算机生成的含 30 个四元正态随机向量的两个样本所作的 $d^{2}$ 图. 正如所期望的, 该图为直线模式, 但最高的两个或三个有序平方距离的变异性很大.

下面给出一个真实的数据集, 可与生成\hyperref[fig:4.8]{图 4.8} 的模拟数据集进行比较.

\noindent\textbf{例 4.14}\quad \textbf{(评估一个四元数据集的多元正态性)}

\hyperref[tab:4.3]{表 4.3} 中的数据是对 $n=30$ 块木板的每一块的硬度的四种不同的测量 $x_{1},x_{2},x_{3}$ 和 $x_{4}$. 第一个测量需要给木板传送一个向下的冲击波, 第二个测量是当木板振动时测得的, 后两个测量是从静态试验中得到的, 平方距离 $d_{j}^{2}=(\mathbf{x}_{j}-\bar{\mathbf{x}})'\mathbf{S}^{-1}(\mathbf{x}_{j}-\bar{\mathbf{x}})$ 也在表中给出.
